% =====================================================================================================
\section{Introduction}\label{sec:intro}
Cameras on doors, corridors, production lines and drones now send video over Wi-Fi and shared Ethernet as a matter of course. Two properties matter for such a link. \emph{Confidentiality} keeps a bystander on the network from watching the scene; \emph{integrity} lets the viewer trust that what it shows was captured by the camera and was not edited, replayed or injected on the way. The usual answer, encrypting the transport with TLS, SRTP~\cite{rfc3711} or a VPN, provides both, but it hides the stream from every tool that is not part of the key infrastructure, and the cipher work grows with the full bitrate.

Selective encryption takes a different route. It encrypts only the parts of the compressed representation that carry most of the visual information and leaves the format intact, so standard decoders still parse the result~\cite{massoudi2008}. For baseline JPEG, the sign and the appended amplitude bits of the quantized DCT coefficients can be encrypted without changing any Huffman code length, so the bitstream stays compliant and its size does not grow~\cite{vandroogenbroeck2002,auer2013}. The price is that such schemes are easy to weaken: the positions of the non-zero coefficients stay visible, and an attacker who guesses the hidden values can still see the outline of the scene~\cite{massoudi2008,lin2024,hirose2024}. Most published schemes also use fixed key material, so one frame the attacker knows exposes every later frame, and almost none authenticate what they encrypt.

This work starts from a practical question: \emph{how much of a secure camera link can be put into the encoder hardware itself, and what does an eavesdropper still learn?} Answering it needs three things that are usually studied apart. The hardware must encrypt at line rate without a frame buffer and without slowing the encoder. The cryptography must be standard, authenticate every frame together with its metadata, and never reuse a nonce across frames, rekeys or reboots. And the security claim must be measured the way an attacker would attack: with replacement attacks, non-zero counting, optimization, puzzle solvers, known plaintext and trained networks, on real video, and against the published schemes under the same conditions.

\subsection{Contributions}
\begin{enumerate}
\item \textit{In-encoder authenticated selective encryption.} A synthesizable Verilog MJPEG encoder in which a keystream of Ascon-AEAD128, the NIST lightweight standard~\cite{sp800232,ascon}, is XORed onto coefficient signs and amplitudes inside the entropy coder (Section~\ref{sec:design}). We prove that the operation preserves every code length (Lemma~\ref{lem:format}) and that the cipher can never stall the encoder (Proposition~\ref{prop:supply}); RTL simulation confirms zero added cycles. A per-frame tag binds the coded scan to the protection level, session, frame counter and image size.
\item \textit{A hardware form of a keyed macroblock permutation.} The strongest level, E5, hides block positions. The fabric codes one restart segment per macroblock with its own keystream and authenticates them in raster order; the ARM core sends them in a keyed per-frame order; the receiver restores the order before the tag check (Section~\ref{sec:e5}). No frame buffer is needed, the eavesdropper's picture is identical to that of a frame-level shuffle, and the order is authenticated.
\item \textit{A parallel-lane architecture.} Restart segments make lanes independent; one, two and four lanes produce identical bytes and reach \RtlOneLaneMpix, \RtlTwoLanesMpix{} and \RtlFourLanesMpix\,Mpixel/s at 100\,MHz in cycle-accurate RTL simulation, close to the analytical bounds of Section~\ref{sec:lanes}. They meet 150\,MHz after place and route on the xc7z020, the complete four-lane PYNQ-Z2 system meets \OverlayLanesFourMhz\,MHz, and on the board every image is bit-exact with the golden model in every level, with encoding clocks within \BdBusyVsRtlMaxPct\,\% of the RTL simulation (Section~\ref{sec:hwresults}).
\item \textit{SVE-Bench, a common yardstick.} Six protection levels and four published format-compliant schemes~\cite{vandroogenbroeck2002,kishore2015,auer2013,hirose2024} are scored under the same four keyless attacks on \BenchSources{} sources, together with jigsaw solvers, sign-bit recovery~\cite{lin2024}, a chosen-plaintext attack~\cite{li2025shuffle}, trained U-Nets and a person detector (Section~\ref{sec:security}).
\item \textit{A closed feedback loop.} An authenticated feedback channel, a forward-secure key ratchet and an adaptive policy raise protection when the receiver reports attacks; we evaluate \NumAttacks{} active attacks and benign impairments, and compare a threshold rule with a contextual bandit over \Seeds{} paired seeds (Section~\ref{sec:protocol}).
\item \textit{A verification flow that keeps hardware and claims honest.} A Python golden model is the referee for every RTL stage; Ascon passes all official known-answer tests; every number in this paper is generated from stored results by a script (Section~\ref{sec:verification}).
\end{enumerate}

\subsection{Organization}
Section~\ref{sec:related} reviews related work and how hardware designs in the field are evaluated. Section~\ref{sec:background} gives the mathematics of baseline JPEG, Ascon and the leakage metrics. Section~\ref{sec:threat} states the threat model. Section~\ref{sec:design} derives the cryptographic design, Section~\ref{sec:hardware} the hardware architecture down to gate level, and Section~\ref{sec:verification} the verification. Sections~\ref{sec:hwresults} to~\ref{sec:further} report hardware, security, protocol and further results; Section~\ref{sec:discussion} discusses limitations and concludes.
